\subsection{Uniform Convergence and VC Theory}

Generalization asks when small training loss implies small population loss.  Uniform convergence controls the gap for all predictors in a class at once.

\begin{definition} \label{def:sl-uniform-convergence} \hlt{Uniform Convergence}\\
A class $\mathcal{F}$ has \hlt{uniform convergence} for a bounded loss if
\[
\sup_{f\in\mathcal{F}}\abs{R(f)-\widehat{R}_S(f)}\to0
\]
in probability, or almost surely, as the sample size tends to infinity.
\end{definition}

\begin{definition} \label{def:sl-vc-dimension} \hlt{VC Dimension}\\
For a class $\mathcal{C}$ of subsets of $\mathcal{X}$, a finite set $\{x_1,\ldots,x_m\}$ is \hlt{shattered} if every subset of it can be realized as $\{x_i:x_i\in C\}$ for some $C\in\mathcal{C}$.  The \hlt{VC dimension} of $\mathcal{C}$ is the largest size of a shattered set, or infinity if no largest finite size exists.
\end{definition}

\begin{theorem} \label{thm:sl-vc-generalization-form} \hlt{VC Generalization Bound Form}\\
For a binary classification class with VC dimension $d$ and zero-one loss, bounds of the form
\[
\sup_{f\in\mathcal{F}}\abs{R(f)-\widehat{R}_S(f)}
\lesssim
\sqrt{\frac{d\log(n/d)+\log(1/\delta)}{n}}
\]
hold with probability at least $1-\delta$ when $n$ is larger than $d$.
\end{theorem}

\begin{remark} \label{rem:sl-growth-functions} \hlt{Growth Functions}\\
The VC dimension controls how many labelings the class can induce on finite samples.  Sauer's lemma converts that combinatorial count into a probability bound.
\end{remark}
